\documentclass[a4paper,12pt]{article}
\usepackage[utf8]{inputenc}
\usepackage[T1]{fontenc}
\usepackage[spanish]{babel}
\usepackage{geometry}
\usepackage{graphicx}
\usepackage{xcolor}
\usepackage{hyperref}
\usepackage{fancyhdr}
\usepackage{enumitem}
\usepackage{amsmath}
\usepackage{booktabs}
\usepackage{array}
\usepackage{setspace}
\usepackage{float}
\usepackage{longtable}
\usepackage{tabularx}
\usepackage{ragged2e}

\geometry{top=34mm,bottom=25mm,left=25mm,right=25mm}

\definecolor{ExternadoGreen}{RGB}{0,104,55}
\definecolor{ExternadoDark}{RGB}{0,51,25}
\definecolor{ExternadoGold}{RGB}{198,146,20}

\hypersetup{colorlinks=true,linkcolor=ExternadoGreen,urlcolor=ExternadoGreen,citecolor=ExternadoGreen}

\setlength{\headheight}{52pt}
\pagestyle{fancy}
\fancyhf{}
\fancyhead[L]{\IfFileExists{firma.jpg}{\includegraphics[height=40pt]{firma.jpg}}{}}
\fancyhead[R]{\textcolor{ExternadoGreen}{\small Departamento de Matemáticas}}
\fancyfoot[C]{\textcolor{ExternadoDark}{\thepage}}

\fancypagestyle{plain}{
  \fancyhf{}
  \fancyhead[L]{\IfFileExists{firma.jpg}{\includegraphics[height=40pt]{firma.jpg}}{}}
  \fancyhead[R]{\textcolor{ExternadoGreen}{\small Universidad Externado de Colombia}}
  \fancyfoot[C]{\textcolor{ExternadoDark}{\thepage}}
}

\title{
\vspace{-8mm}
\textbf{\textcolor{ExternadoGreen}{Ficha de Proyecto de Investigación}}\\
\large \textcolor{ExternadoDark}{Grupo de Investigación en Ciencia de Datos}\\
\large \textcolor{ExternadoDark}{Universidad Externado de Colombia}
}
\author{\textcolor{ExternadoDark}{Departamento de Matemáticas}}
\date{\textcolor{ExternadoDark}{\today}}

\setlist[itemize]{leftmargin=*, itemsep=3pt, topsep=3pt}
\setstretch{1.12}

\begin{document}
\maketitle
\vspace{-3mm}
\noindent\textcolor{ExternadoDark}{\rule{\textwidth}{0.6pt}}

\section*{1. Información general}
\begin{longtable}{p{4.2cm} p{10.4cm}}
\toprule
\textbf{Campo} & \textbf{Detalle} \\
\midrule
Título del proyecto & Utilización del policarbonato alveolar para colectores termosolares y como purificador de agua \\
Investigador principal & Leonardo Guarín Salomón \\
Co-investigadores & No reportado \\
Líneas de investigación & Modelado Matemático \\
Año de inicio & 2025-07-01 \\
Duración estimada & No reportado. \\
Estado del proyecto & Ejecución \\
Tipo de proyecto & Investigación académica;Desarrollo aplicado; \\
Nivel de avance actual & Idea inicial;Diseño metodológico;Recolección de datos;Desarrollo / implementación; \\
Semillero vinculado & No \\
Nombre del semillero & No aplica \\
Fuentes de financiación & Empresas privadas; \\
\bottomrule
\end{longtable}

\section*{2. Problema de investigación}
El calentamiento global amenaza la existencia de la humanidad y es un deber combatirlo, las emisiones de gases de efecto invernadero deben reducirse con urgencia pues nos estamos acercando a escenarios catastróficos, la Universidad tiene el deber de servir de guía hacia la solución del problema, planteando soluciones viables que ayuden al logro de los Objetivos de Desarrollo Sostenible (ODS) al disminuir las emisiones de gases de efecto invernadero y producir soluciones viables para poblaciones que están consumiendo agua con patógenos.

\section*{3. Objetivo general}
Es en este escenario donde nos planteamos las preguntas: 1. ¿Podemos producir, con materiales y medios locales, un calentador termosolar que sea competitivo en la relación duración/precio con los calentadores de tubos al vacío producidos en China? Cabe anotar que esta relación es actualmente de 2.4 años por millón de pesos y estimamos que podemos llevarla a un mínimo de 6 años por millón de pesos. Se busca que cualquier persona en cualquier pueblo o ciudad pueda armar un calentador solar plenamente funcional. Para la solución de este problema proponemos la utilización de láminas de policarbonato alveolar como superficie colectora, logrando con su utilización un abaratamiento significativo en el costo final del calentador solar. 2. ¿Constituyen las láminas de policarbonato alveolar una solución para eliminar patógenos del agua? Se busca construir un sistema de purificación de agua, económico y eficiente que sirva para poblaciones que no tengan acceso a agua tratada.

\section*{4. Metodología}
Este es un proceso de experimentación con diferentes configuraciones de materiales para la construcción de calentadores y purificadores de agua con energía solar. Los materiales se consiguen a través de donaciones.

\section*{5. Comunidades a impactar}
\subsection*{5.1. Comunidades externas}
\begin{itemize}
    \item Comunidades del sector rural.
\end{itemize}

\subsection*{5.2. Comunidades internas}
\begin{itemize}
    \item Los estudiantes a través del curso Arme su calentador Solar.
\end{itemize}

\section*{6. Semillero y articulación formativa}
\subsection*{6.1. Participación del semillero}
\begin{itemize}
    \item ninguna
\end{itemize}

\subsection*{6.2. Aliados externos}
No reportado.

\section*{7. Requerimientos tecnológicos}
\begin{itemize}
    \item Elementos electrónicos para recolectar los datos del calentador.
\end{itemize}

\section*{8. Resultados esperados}
\subsection*{8.1. Resultados generales}
1. Se espera producir un calentador solar con indicadores de costo, eficiencia y duración competitivos en el mercado. 2. Se espera producir un purificador de agua económico y con indicadores de costo inferiores a los del mercado.

\subsection*{8.2. Resultados científicos esperados}
\begin{itemize}
    \item Artículo científico
    \item Working paper
    \item Ponencia
\end{itemize}

\subsection*{8.3. Productos aplicados}
\begin{itemize}
    \item Prototipo
\end{itemize}

\section*{9. Observación de gestión}
Este documento fue generado automáticamente a partir del formulario de registro de proyectos del grupo. Se recomienda revisar y ajustar la redacción final antes de su circulación institucional o envío a convocatorias.

\end{document}
